\documentclass{ximera}
\title{Definite intgral}
\begin{document}
\begin{abstract}
Partitions of an interval, Riemann sums (right-hand, left-hand and midpoint), the definition of the definite integral as a limit of Riemann sums, some history of the integral sign, and a first integral computed from the definition. Textbook §5.1–5.2.
\end{abstract}
\maketitle

Now let's look at a different function, and we're going to be a little more precise in how we are defining these rectangles. We will work with $f(x)=x^{5}-6 x^{2}-$ $x+1$ :
\begin{image}
\includegraphics[width=0.6\textwidth]{b48a5324-93dc-45ee-b7c2-dc0b5706bda6-135_563_769_447_670}
\end{image}

Given a closed interval $[a, b]$, a partition of an interval of numbers $x_{i}$ is such that:
\[
a=x_{0}<x_{1}<x_{2}<\ldots<x_{n-1}<x_{n}=b .
\]

This gives us little sub-intervals:
\[
\left[x_{0}, x_{1}\right],\left[x_{1}, x_{2}\right], \ldots,\left[x_{n-1}, x_{n}\right]
\]

If the intervals are all the same size $\left(\frac{b-a}{n}\right)$ we say that the partition is regular.

\begin{example}\label{example:12-7}
For some examples, let $a=2$ and $b=5$. The following are two different partitions:
\[
\begin{gathered}
2<2.1<3<4.2<4.5<4.51<5 \\
2<3<4<5
\end{gathered}
\]

and they each define different intervals:
\[
\begin{gathered}
{[2,2.1],[2.1,3],[3,4.2],[4.2,4.5],[4.5,4.51],[4.51,5]} \\
{[2,3],[3,4],[4,5]}
\end{gathered}
\]

Which partition is regular? \wordChoice{\choice[correct]{the second, $2<3<4<5$}\choice{the first, $2<2.1<3<4.2<4.5<4.51<5$}\choice{both}\choice{neither}}
\end{example}

\begin{exercise}\label{exercise:12-8}
Given $a=0$ and $b=1$, give a regular partition:
\begin{freeResponse}\end{freeResponse}
\end{exercise}

Always assume regular The length of an interval is the difference between the start and end of the interval: $\Delta x_{i}=\answer{x_{i}-x_{i-1}}$ for the interval $\left[x_{i-1}, x_{i}\right]$.

We are now ready to define our summations! The Riemann sum of a function on an interval $[a, b]$ is the sum:
\[
\answer{\sum_{i=1}^{n} f\left(x_{i}^{\star}\right) \Delta x_{i}}
\]
where $x_{i}^{\star} \in\left[x_{i-1}, x_{i}\right]$ is some arbitrary number.
\begin{image}
\includegraphics[width=0.45\textwidth]{b48a5324-93dc-45ee-b7c2-dc0b5706bda6-136_839_570_1222_489}
\end{image}
\begin{image}
\includegraphics[width=0.45\textwidth]{b48a5324-93dc-45ee-b7c2-dc0b5706bda6-136_837_570_1222_1061}
\end{image}
% blank grids for sketching Riemann sums in the original; open work per spec 3(e)
\begin{freeResponse}\end{freeResponse}

We've already seen two of these! The right-hand sum is the Riemann sum where we let $x_{i}^{\star}=x_{i}$. The left-hand sum is the Riemann sum where we let $x_{i}^{\star}=x_{i-1}$. The midpoint sum is the Riemann sum where we let $x_{i}^{\star}=\frac{x_{i-1}+x_{i}}{2}$.

But what we reeeeaaaallllly want is to find the exact area, not just approximations! What can we do to find the exact area?

We're going to use our favorite tool so far in this course: limits. Yes. . . limits.

The limit we're going to use is actually the following limit:
\[
\answer{\lim_{n \rightarrow \infty} \sum_{i=1}^{n} f\left(x_{i}^{\star}\right) \Delta x_{i}}
\]
If this limit exists we call say that the function is Riemann-integrable.

BUT, this notation is horrible. There are a million and a half symbols. Who wants to write that each time? No one. So we're going to use another notation that (I think) everyone has already seen: the integral.

\begin{remark}\label{remark:12-9}
Some history! We use $\Sigma$ for summation because summation starts with s. So, we need a new symbol for the integral. Luckily, Gottfried Leibniz (German) had came up with the perfect notation back in 1675! He used ancient German's long s: ſ. (Nowadays, German's long s is written as "ß".) Why? Because the integral is a summation of realy small intervals. So ſ is perfect for that since it keeps it as an "s", but is a different s! (Even in English we used to have ſ! It wasn't until between 1800 and 1820 that this letter disappeared from English.)
\end{remark}

We changed the "s", but we also need to change the $\Delta$ since we're taking the limit. Since we're taking the limit, these intervals are getting smaller and smaller. In other words, it would be nice to represent that by using a "small" $\Delta$. So what's the lower-case of $\Delta$ ? It's $\delta$ of course! But, Leibniz, who was one of the founders of calculus, translated the greek to German and the $\delta$ became a $d$. This is where Leibniz notation: $\frac{d y}{d x}$ actually came from!

So, if our function is Riemann-integrable, we will write:
\[
\answer{\int_{a}^{b} f(x)\,dx} = \lim _{n \rightarrow \infty} \sum_{i=1}^{n} f\left(x_{i}^{\star}\right) \Delta x_{i} .
\]

Some definitions:
\begin{itemize}
  \item The $a$ is the lower limit of integration,
  \item The $b$ is the upper limit of integration,
  \item The $f(x)$ is the integrand,
  \item The $x$ is called the variable of integration or the independent variable.
\end{itemize}

\begin{theorem}\label{theorem:12-10}
% CHECKED-FIX: was distractors {differentiable}{increasing}{bounded} -- "differentiable" and "increasing" also give true statements (both imply integrable), so they were not wrong
If a function $f(x)$ is \wordChoice{\choice[correct]{continuous}\choice{bounded}\choice{defined}\choice{positive}} on the interval $[a, b]$ or if $f(x)$ has only a finite number of jump discontinuities, then the function $f(x)$ is \wordChoice{\choice[correct]{Riemann-integrable}\choice{differentiable}\choice{continuous}\choice{constant}} on $[a, b]$.
\end{theorem}

\begin{example}\label{example:12-11}
Let's look at an (easy) example. Find $\int_{1}^{2} x^{2} d x$.
% work: regular partition of [1,2] into n pieces, right-hand endpoints: Delta x = 1/n, x_i = 1 + i/n;
% sum_{i=1}^{n} (1 + i/n)^2 (1/n) = 1 + (2/n^2) n(n+1)/2 + (1/n^3) n(n+1)(2n+1)/6 -> 1 + 1 + 1/3 = 7/3
\[
\Delta x=\answer{\frac{1}{n}}, \qquad x_{i}=\answer{1+\frac{i}{n}}
\]
\[
\int_{1}^{2} x^{2} d x=\lim _{n \rightarrow \infty} \sum_{i=1}^{n} f\left(x_{i}\right) \Delta x=\answer{\frac{7}{3}}
\]
\end{example}

That was disgusting. The whole point of the next few days (and the next class) is to try and make this easier.

\end{document}
